\section{Saturation of the area law}

\begin{definition}[Block entropy of an MPDO]\label{def:block_entropy}
    \lean{MPOTensor.blockEntropy}
    \leanok
    \uses{def:mpo_tensor, def:von_neumann_entropy}
    For a fixed system size $N$, a block length $L\le N$, and an MPO tensor $M$
    such that $\rho^{(N)}(M)$ is positive semidefinite, the \emph{$L$-block entropy}
    $S_L^{(N)}(M)$ is the von Neumann entropy of the reduced state of the first
    $L$ of $N$ spins of the normalized state
    $\sigma^{(N)}(M) = \rho^{(N)}(M) / \tr[\rho^{(N)}(M)]$.
\end{definition}

\begin{definition}[Mutual information of an MPDO]\label{def:mutual_info_chain}
    \lean{MPOTensor.mutualInfoChain}
    \leanok
    \uses{def:block_entropy}
    For the same fixed system size $N$, block length $L\le N$, and positive
    semidefiniteness of $\rho^{(N)}(M)$, the \emph{mutual information} between
    a block of $L$ spins and the rest is
    $I_L = S_L^{(N)}(M) + S_{N-L}^{(N)}(M) - S_N^{(N)}(M)$.
\end{definition}

\begin{definition}[Saturation of the area law]\label{def:is_sal}
    \lean{MPOTensor.IsSAL}
    \leanok
    \uses{def:mutual_info_chain}
    An MPO tensor $M$ that generates an MPDO and whose finite-chain operators
    $\rho^{(N)}(M)$ all have nonzero trace (so the normalized states are
    well-defined) \emph{saturates the area law} if
    $I_L = I_{L+1}$ for all $1 \le L < \lfloor N/2 \rfloor$ and all $N$; that is,
    the mutual information is constant, $I_1 = I_2 = \cdots = I_{\lfloor N/2 \rfloor}$.
    This is \cite[Definition~4.6, line~811]{Cirac2016MPDO_arXiv}.
\end{definition}

\begin{lemma}[Saturation telescopes to a constant mutual information]
    \label{lem:sal_const}
    \lean{MPOTensor.mutualInfoChain_eq_of_isSAL}
    \leanok
    \uses{def:is_sal}
    If an MPO tensor $M$ saturates the area law, then all mutual informations in
    the range $1\le L\le\lfloor N/2\rfloor$ coincide:
    \[
        I_L = I_{L'}
        \qquad
        \text{for all } 1\le L, L'\le \lfloor N/2\rfloor .
    \]
\end{lemma}
\begin{proof}\leanok
    \uses{def:is_sal}
    The defining condition supplies the single-step equalities $I_m = I_{m+1}$ for
    $1\le m<\lfloor N/2\rfloor$. Climbing one step at a time from the smaller of $L$
    and $L'$ up to the larger composes these into the stated equality.
\end{proof}

\begin{proposition}[Monotonicity of the mutual information]
    \label{prop:mutual_info_monotone}
    \lean{mutualInfoChain_monotone}
    \leanok
    \uses{def:mutual_info_chain}
    Let $M$ generate an MPDO whose finite-chain operator $\rho^{(N)}(M)$ is
    positive semidefinite with nonzero trace. For every block length $L$ with
    $2L+1\le N$,
    \[
        I_L \le I_{L+1}.
    \]
    The implication
    \[
        1\le L<\lfloor N/2\rfloor \quad\Longrightarrow\quad 2L+1\le N
    \]
    shows that this contains the range of
    \cite[Proposition~4.5, line~801]{Cirac2016MPDO_arXiv}.
\end{proposition}
\begin{proof}\leanok
    \uses{thm:entropy_strong_subadditivity, def:block_entropy}
    Apply strong subadditivity to the reduced state of the first $N-L$ spins,
    split into three contiguous segments of lengths $1$, $L$, and $N-2L-1$.
    Tracing the appropriate segments yields the four block entropies
    $S_{N-L}$, $S_L$, $S_{L+1}$, and $S_{N-L-1}$, where the cyclic invariance of
    $\rho^{(N)}(M)$ identifies the entropy of a contiguous block with that of any
    block of the same length. Strong subadditivity then gives
    $S_{N-L} + S_L \le S_{L+1} + S_{N-L-1}$, which rearranges to
    $I_L \le I_{L+1}$ after subtracting $S_N$.
\end{proof}

\begin{lemma}[Nonnegativity of block entropies]
    \label{lem:block_entropy_nonneg}
    \lean{blockEntropy_nonneg}
    \leanok
    \uses{def:block_entropy, thm:entropy_nonneg}
    For a normalized positive semidefinite finite-chain operator, every block entropy
    is nonnegative: $0 \le S_m$.
\end{lemma}
\begin{proof}\leanok
    The reduced state of a block is a density matrix, and the von Neumann entropy of
    a density matrix is nonnegative.
\end{proof}

\begin{lemma}[Nonnegativity of the mutual information]
    \label{lem:mutual_info_chain_nonneg}
    \lean{mutualInfoChain_nonneg}
    \leanok
    \uses{def:mutual_info_chain, thm:entropy_strong_subadditivity, lem:block_entropy_nonneg}
    Let $M$ generate an MPDO whose finite-chain operator $\rho^{(N)}(M)$ is positive
    semidefinite with nonzero trace. For every block length $L\le N$,
    \[
        0 \le I_L.
    \]
\end{lemma}
\begin{proof}\leanok
    \uses{thm:entropy_strong_subadditivity, def:block_entropy, lem:block_entropy_nonneg}
    Apply strong subadditivity with a trivial middle subsystem (equivalently,
    subadditivity) to the bipartition of the chain into the first $L$ and last
    $N-L$ spins: $S_N + S_0 \le S_L + S_{N-L}$. Since $S_0 \ge 0$
    (Lemma~\ref{lem:block_entropy_nonneg}), rearranging gives
    $I_L = S_L + S_{N-L} - S_N \ge S_0 \ge 0$.
\end{proof}

\begin{lemma}[Symmetry of the mutual information across the cut]
    \label{lem:mutual_info_chain_symm}
    \lean{mutualInfoChain_symm}
    \leanok
    \uses{def:mutual_info_chain}
    For every block length $L\le N$, the mutual information of the first $L$ spins
    equals that of their complement:
    \[
        I_L = I_{N-L}.
    \]
\end{lemma}
\begin{proof}\leanok
    \uses{def:mutual_info_chain}
    Both sides equal $S_L + S_{N-L} - S_N$ by the symmetric definition
    $I_L = S_L + S_{N-L} - S_N$, using $N-(N-L) = L$.
\end{proof}

\begin{lemma}[Empty-block entropy vanishes]
    \label{lem:block_entropy_zero}
    \lean{blockEntropy_zero}
    \leanok
    \uses{def:block_entropy, thm:entropy_le_log_rank, lem:block_entropy_nonneg}
    For a normalized positive semidefinite finite-chain operator, the block entropy of
    the empty block vanishes: $S_0 = 0$.
\end{lemma}
\begin{proof}\leanok
    \uses{thm:entropy_le_log_rank, lem:block_entropy_nonneg}
    The reduced state of zero spins is one-dimensional, so its rank is one and the rank
    bound on the entropy gives $S_0 \le \log 1 = 0$; with $S_0 \ge 0$ this forces
    $S_0 = 0$.
\end{proof}

\begin{lemma}[Empty-block mutual information vanishes]
    \label{lem:mutual_info_chain_zero}
    \lean{mutualInfoChain_zero}
    \leanok
    \uses{def:mutual_info_chain, lem:block_entropy_zero}
    For a normalized positive semidefinite finite-chain operator, the mutual information
    of the empty block vanishes: $I_0 = 0$.
\end{lemma}
\begin{proof}\leanok
    \uses{lem:block_entropy_zero}
    By definition $I_0 = S_0 + S_N - S_N = S_0$, which is zero by
    Lemma~\ref{lem:block_entropy_zero}.
\end{proof}

\begin{remark}[Boundedness of the mutual information]
    The monotone sequence $I_L$ is bounded by a constant depending only on the
    bond dimension, $I_L \le 4\log D$, so it converges as $L\to\infty$. For a
    general mixed MPDO this bound is not proved in
    \cite{Cirac2016MPDO_arXiv}; that reference attributes it (line~1319) to the
    mutual-information area law of Wolf et al., used there as an external result.
    The pure-state instance is formalized in
    Proposition~\ref{prop:pure_mutual_info_bound} below.
\end{remark}

\begin{definition}[Block entropy and area-law saturation for pure states]
    \label{def:pure_sal}
    \lean{MPSTensor.pureBlockEntropy, MPSTensor.IsSAL}
    \leanok
    \uses{def:mps_tensor, def:von_neumann_entropy}
    For an MPS tensor $A$, a system size $N$, and a block length $L\le N$,
    the $L$-block entropy $S_L^{(N)}(A)$ is the von Neumann entropy of the
    reduced state of the first $L$ spins of
    \[
        \sigma^{(N)}(A)
        =
        \frac{\ketbra{V^{(N)}(A)}{V^{(N)}(A)}}{\|V^{(N)}(A)\|^2}.
    \]
    The tensor saturates the area law if
    $S_L^{(N)}(A) = S_{L+1}^{(N)}(A)$ for all
    $1 \le L < \lfloor N/2 \rfloor$ and all $N$.
    This is \cite[Definition~3.13, line~600]{Cirac2016MPDO_arXiv}.
\end{definition}

\begin{definition}[Purification tensor]\label{def:doubled_tensor}
    \lean{doubledTensor}
    \leanok
    \uses{def:mps_tensor, def:mpo_tensor}
    For an MPS tensor $A$ on bond space $\mathbb{C}^D$, define the associated
    MPO tensor $M$ on bond space $\mathbb{C}^{D^2}$ by
    \[
        M^{ij} = A^i \otimes \overline{A^j}.
    \]
    This is the single-ancilla purification tensor associated with
    \cite[eq.~(line~752)]{Cirac2016MPDO_arXiv}.
\end{definition}

\begin{lemma}[Pure state as a purification MPDO]\label{lem:mpo_doubled}
    \lean{mpo_doubledTensor}
    \leanok
    \uses{def:mps_tensor, def:mpo_family, def:doubled_tensor}
    For an MPS tensor $A$, the pure-state operator
    $\ketbra{V^{(N)}(A)}{V^{(N)}(A)}$ equals the MPDO generated by the tensor
    $M^{ij} = A^i \otimes \overline{A^j}$, acting on a bond space of
    dimension $D^2$:
    \[
        \rho^{(N)}\bigl(A \otimes \bar A\bigr) = \ketbra{V^{(N)}(A)}{V^{(N)}(A)}.
    \]
    This is the single-ancilla case of the purification picture of
    \cite[eq.~(line~752)]{Cirac2016MPDO_arXiv}; it lets the MPDO block-entropy
    theory apply to pure-state block entropies.
\end{lemma}
\begin{proof}\leanok
    The word product of doubled-tensor entries is the Kronecker product of the
    bra word product and the conjugated ket word product (mixed-product property
    of the Kronecker product), so its trace factors as
    $\overline{V^{(N)}(A)(\tau)}\,V^{(N)}(A)(\sigma)$, the matrix element of
    $\ketbra{V^{(N)}(A)}{V^{(N)}(A)}$.
\end{proof}

\begin{lemma}[Block entropy of the purification tensor]
    \label{lem:blockEntropy_doubledTensor}
    \lean{blockEntropy_doubledTensor}
    \leanok
    \uses{lem:mpo_doubled, def:block_entropy, def:pure_sal}
    For an MPS tensor $A$, system size $N$, and block length $L \le N$, the MPDO
    block entropy of the tensor $M^{ij} = A^i \otimes \overline{A^j}$ agrees
    with the pure-state block entropy:
    \[
        S_L^{(N)}(M) = S_L^{(N)}(A).
    \]
\end{lemma}
\begin{proof}\leanok
    Lemma~\ref{lem:mpo_doubled} gives
    $\rho^{(N)}(M) = \ketbra{V^{(N)}(A)}{V^{(N)}(A)}$, hence
    \[
        \sigma^{(N)}(M)
        =
        \frac{\rho^{(N)}(M)}{\tr[\rho^{(N)}(M)]}
        =
        \frac{\ketbra{V^{(N)}(A)}{V^{(N)}(A)}}{\|V^{(N)}(A)\|^2}
        =
        \sigma^{(N)}(A).
    \]
    Therefore the reduced block states agree:
    \[
        \rho_L^{(N)}(M)
        =
        \tr_{N\setminus L}\!\bigl[\sigma^{(N)}(M)\bigr]
        =
        \tr_{N\setminus L}\!\bigl[\sigma^{(N)}(A)\bigr]
        =
        \rho_L^{(N)}(A).
    \]
    Thus
    \[
        S_L^{(N)}(M)
        =
        S\bigl(\rho_L^{(N)}(M)\bigr)
        =
        S\bigl(\rho_L^{(N)}(A)\bigr)
        =
        S_L^{(N)}(A).
    \]
\end{proof}

\begin{lemma}[Schmidt symmetry of the block entropy]\label{lem:schmidt_symmetry}
    \lean{pureBlockEntropy_complement}
    \leanok
    \uses{def:pure_sal}
    For an MPS tensor $A$ and a block length $L\le N$, the block entropies of a
    block and its complement coincide:
    \[
        S_L^{(N)}(A) = S_{N-L}^{(N)}(A).
    \]
\end{lemma}
\begin{proof}\leanok
    The reduced state of the first $L$ spins is $\rho_L \propto W W^\dagger$, where
    $W$ is the matrix of amplitudes $\langle u\,w\,|\,V^{(N)}(A)\rangle$ with $u$ the
    first $L$ spins and $w$ the remaining $N-L$. By cyclic invariance of the
    amplitudes, the reduced state of the complement is $\propto \overline{W^\dagger W}$.
    Since $W W^\dagger$ and $W^\dagger W$ share the same nonzero spectrum, and
    complex conjugation preserves the (real) eigenvalues of a Hermitian matrix, the
    two reduced states have equal von Neumann entropy.
\end{proof}

\begin{proposition}[Pure-state area-law monotonicity]\label{prop:pure_area_monotone}
    \lean{pureBlockEntropy_monotone}
    \leanok
    \uses{def:pure_sal, lem:schmidt_symmetry, lem:blockEntropy_doubledTensor}
    Let $A$ be an MPS tensor with $V^{(N)}(A)\neq 0$. For every block length $L$
    with $2L+1\le N$ (covering $1\le L<\lfloor N/2\rfloor$),
    \[
        S_L^{(N)}(A) \le S_{L+1}^{(N)}(A).
    \]
    This is the pure-state area law of \cite[Section~3, line~599]{Cirac2016MPDO_arXiv}.
\end{proposition}
\begin{proof}\leanok
    \uses{prop:mutual_info_monotone}
    Apply the block-entropy form of strong subadditivity (the inequality
    underlying Proposition~\ref{prop:mutual_info_monotone}) to the purification
    MPDO of $A$, giving $S_{N-L} + S_L \le S_{L+1} + S_{N-L-1}$. By the Schmidt
    symmetry $S_{N-L} = S_L$ and $S_{N-L-1} = S_{L+1}$, this reduces to
    $S_L \le S_{L+1}$.
\end{proof}

\begin{lemma}[Operator-Schmidt Gram is a transfer-map power]
    \label{lem:schmidt_gram_transfer_power}
    \lean{MPSTensor.schmidtLeft, MPSTensor.schmidtRight,
        MPSTensor.schmidtLeft_gram_apply, MPSTensor.schmidtRight_gram_apply}
    \leanok
    \uses{def:transfer_map}
    Let $A$ be an MPS tensor of bond dimension $D$. For a length-$\ell$
    configuration $\sigma$ and a length-$m$ configuration $\tau$, write the
    operator-Schmidt factors as
    \[
        (L_\ell)_{\sigma,(a_1,a_2)} = (A^\sigma)_{a_1a_2},
        \qquad
        (R_m)_{(p_1,p_2),\tau} = (A^\tau)_{p_2p_1}.
    \]
    The $D^2\times D^2$ Gram matrix of the left factor equals the $\ell$-fold
    transfer map evaluated on a matrix unit:
    \[
        \bigl(L_\ell^\dagger L_\ell\bigr)_{(a_1,a_2),(b_1,b_2)}
            = \sum_{\sigma}\overline{(A^\sigma)_{a_1 a_2}}\,(A^\sigma)_{b_1 b_2}
            = \bigl(\mathbb{E}_A^{\,\ell}(\ketbra{b_2}{a_2})\bigr)_{b_1,a_1},
    \]
    the sum running over length-$\ell$ configurations $\sigma$. The Gram of the
    right factor $R_m$ of an $m$-spin complement satisfies the bond-swapped
    identity
    \[
        \bigl(R_m R_m^\dagger\bigr)_{(p_1,p_2),(q_1,q_2)}
            = \sum_{\tau}(A^\tau)_{p_2 p_1}\,
                \overline{(A^\tau)_{q_2 q_1}}
            = \bigl(\mathbb{E}_A^{\,m}(\ketbra{p_1}{q_1})\bigr)_{p_2,q_2},
    \]
    the sum running over length-$m$ configurations $\tau$.
\end{lemma}
\begin{proof}\leanok
    The left Gram entry is
    $\sum_\sigma \overline{(A^\sigma)_{a_1 a_2}}\,(A^\sigma)_{b_1 b_2}$
    by definition of $L_\ell$. Expanding
    $\mathbb{E}_A^{\,\ell}(\ketbra{b_2}{a_2})
        = \sum_\sigma A^\sigma\,\ketbra{b_2}{a_2}\,(A^\sigma)^\dagger$
    and reading off the $(b_1,a_1)$ entry gives the same sum, since
    $\bigl(A^\sigma\,\ketbra{b_2}{a_2}\,(A^\sigma)^\dagger\bigr)_{b_1,a_1}
        = (A^\sigma)_{b_1 b_2}\,\overline{(A^\sigma)_{a_1 a_2}}$.
    The right-factor formula is the same expansion with the two boundary bonds
    interchanged: the $(p,q)$ entry of $R_mR_m^\dagger$ is
    $\sum_\tau (A^\tau)_{p_2p_1}\overline{(A^\tau)_{q_2q_1}}$, which is the
    $(p_2,q_2)$ entry of
    $\mathbb{E}_A^{\,m}(\ketbra{p_1}{q_1})$.
\end{proof}

\begin{lemma}[RFP collapse of the operator-Schmidt Grams]
    \label{lem:rfp_schmidt_gram_collapse}
    \lean{MPSTensor.schmidtLeft_gram_eq_of_isRFP,
        MPSTensor.schmidtRight_gram_eq_of_isRFP}
    \leanok
    \uses{def:is_rfp, lem:schmidt_gram_transfer_power}
    Let $A$ be a renormalization fixed point. For every $L\ge 1$ and
    $m\ge 1$, the operator-Schmidt Grams of the $L$-block and of the
    $m$-site complement are the one-site Grams:
    \[
        L_L^\dagger L_L = L_1^\dagger L_1,
        \qquad
        R_mR_m^\dagger = R_1R_1^\dagger .
    \]
\end{lemma}
\begin{proof}\leanok
    The RFP hypothesis says that the transfer map is idempotent. Hence
    $\mathbb{E}_A^{\,k}=\mathbb{E}_A$ for every $k\ge 1$. Applying
    Lemma~\ref{lem:schmidt_gram_transfer_power} to the left and right Gram
    entries gives the two displayed identities.
\end{proof}

\begin{lemma}[Pure block entropy as a bond-environment charpoly sum]
    \label{lem:pure_block_entropy_env_charpoly}
    \lean{MPSTensor.pureBlockEntropy_eq_env_charpoly}
    \leanok
    \uses{lem:schmidt_gram_transfer_power, lem:charpoly_root_entropy_cyclic_swap}
    Let $A$ be an MPS tensor, let $L\le N$, put $m=N-L$, and write
    $c=(\tr\rho^{(N)})^{-1}$. Then the pure block entropy is the
    charpoly-root entropy sum of the $D^2\times D^2$ bond environment:
    \[
        S_L^{(N)}(A)
        =
        \sum_{\lambda\in
            \operatorname{roots}\chi_{c\,(R_mR_m^\dagger)(L_L^\dagger L_L)}}
            -\operatorname{Re}(\lambda)\log\operatorname{Re}(\lambda),
    \]
    with roots counted with algebraic multiplicity.
\end{lemma}
\begin{proof}\leanok
    The reduced block state has the exact factorization
    \[
        \rho_L^{(N)}(A)
        =
        c\,L_L(R_mR_m^\dagger)L_L^\dagger .
    \]
    Cyclic invariance of the charpoly-root entropy sum gives
    \[
        S\bigl(c\,L_L(R_mR_m^\dagger)L_L^\dagger\bigr)
            =
        S\bigl(c\,(R_mR_m^\dagger)(L_L^\dagger L_L)\bigr),
    \]
    which is the displayed formula.
\end{proof}

\begin{theorem}[A renormalization fixed point saturates the area law]
    \label{thm:rfp-saturates-area-law}
    \lean{MPSTensor.isSAL_of_isRFP}
    \leanok
    \uses{def:is_rfp, def:pure_sal, lem:pure_block_entropy_env_charpoly,
        lem:rfp_schmidt_gram_collapse}
    If an MPS tensor $A$ is a renormalization fixed point, then it saturates the
    area law: for every chain length $N$ and every block length
    $1\le L<\lfloor N/2\rfloor$, the pure-state block entropy is constant in the
    block size,
    \[
        S_L^{(N)}(A) = S_{L+1}^{(N)}(A).
    \]
    The source statement is
    \cite[Proposition~ZCLandSALpure, lines~606--608]{Cirac2016MPDO_arXiv}.
    It assumes canonical form; the theorem above is the stronger pure-state
    implication obtained from transfer-map idempotence alone.
\end{theorem}
\begin{proof}\leanok
    \uses{lem:pure_block_entropy_env_charpoly, lem:rfp_schmidt_gram_collapse}
    Write $c=(\tr\rho^{(N)})^{-1}$. Lemma~\ref{lem:pure_block_entropy_env_charpoly}
    computes the entropy from the bond environment:
    \[
        S_L^{(N)}(A)
        =
        \sum_{\lambda\in
            \operatorname{roots}\chi_{c\,(R_{N-L}R_{N-L}^\dagger)
                (L_L^\dagger L_L)}}
            -\operatorname{Re}(\lambda)\log\operatorname{Re}(\lambda).
    \]
    Since $1\le L<\lfloor N/2\rfloor$, the four lengths
    $L$, $L+1$, $N-L$, and $N-(L+1)$ are positive. The RFP Gram-collapse lemma
    gives
    \[
        L_L^\dagger L_L=L_1^\dagger L_1,\qquad
        L_{L+1}^\dagger L_{L+1}=L_1^\dagger L_1,
    \]
    and
    \[
        R_{N-L}R_{N-L}^\dagger=R_1R_1^\dagger,\qquad
        R_{N-(L+1)}R_{N-(L+1)}^\dagger=R_1R_1^\dagger .
    \]
    Therefore the two bond environments agree:
    \[
        (R_{N-L}R_{N-L}^\dagger)(L_L^\dagger L_L)
        =
        (R_1R_1^\dagger)(L_1^\dagger L_1)
        =
        (R_{N-(L+1)}R_{N-(L+1)}^\dagger)
            (L_{L+1}^\dagger L_{L+1}).
    \]
    The scalar $c$ depends only on $N$, so the two charpoly-root entropy sums
    are equal. Hence $S_L^{(N)}(A)=S_{L+1}^{(N)}(A)$.
\end{proof}

\begin{theorem}[A renormalization fixed point has a constant block-entropy chain]
    \label{thm:rfp_constant_entropy}
    \lean{MPSTensor.pureBlockEntropy_eq_of_isRFP}
    \leanok
    \uses{thm:rfp-saturates-area-law, lem:pure_sal_const}
    If an MPS tensor $A$ is a renormalization fixed point, then all block entropies
    in the range $1\le L\le\lfloor N/2\rfloor$ coincide:
    \[
        S_1^{(N)}(A) = S_2^{(N)}(A) = \cdots = S_{\lfloor N/2\rfloor}^{(N)}(A).
    \]
    This is the explicit constant-entropy chain of the area-law saturation
    \cite[Definition~3.13, line~600]{Cirac2016MPDO_arXiv}, under the fixed-point
    hypothesis.
\end{theorem}
\begin{proof}\leanok
    \uses{thm:rfp-saturates-area-law, lem:pure_sal_const}
    The fixed point saturates the area law
    (Theorem~\ref{thm:rfp-saturates-area-law}), so by the telescoping of
    saturation (Lemma~\ref{lem:pure_sal_const}) the block entropies at any two
    lengths $1\le L, L'\le\lfloor N/2\rfloor$ agree:
    \[
        S_L^{(N)}(A) = S_{L'}^{(N)}(A).
    \]
\end{proof}

\begin{lemma}[Operator-Schmidt rank bound]\label{lem:operator_schmidt_rank}
    \lean{MPSTensor.rank_reducedPureBlockState_le}
    \leanok
    For an MPS tensor $A$ of bond dimension $D$ and a block length $L\le N$, the
    reduced state $\rho_L^{(N)}(A)$ of the first $L$ spins has rank at most $D^2$.
\end{lemma}
\begin{proof}\leanok
    The reduced state is $\rho_L \propto W W^\dagger$, where the amplitude matrix
    $W(u, w) = \langle u\,w\,|\,V^{(N)}(A)\rangle = \operatorname{tr}
        \bigl(A^{u_1}\cdots A^{u_L}\, A^{w_1}\cdots A^{w_{N-L}}\bigr)$ factors through
    the two bond pairs cut by the block boundary. Writing the trace as a sum over
    the bond index pair $(a, b)\in\{1,\dots,D\}^2$ exhibits $W = L'R'$ with $L'$
    having $D^2$ columns, so $\operatorname{rank}\rho_L \le \operatorname{rank} W
        \le D^2$.
\end{proof}

\begin{lemma}[The reduced pure block state is a density matrix]
    \label{lem:reduced_pure_block_state_density}
    \lean{MPSTensor.reducedPureBlockState_posSemidef, MPSTensor.reducedPureBlockState_trace}
    \leanok
    \uses{def:pure_sal}
    For an MPS tensor $A$ with $V^{(N)}(A)\neq 0$, the reduced state $\rho_L^{(N)}(A)$ of
    the first $L$ spins is a density matrix: it is positive semidefinite and has unit
    trace.
\end{lemma}
\begin{proof}\leanok
    Positive semidefiniteness is preserved by the partial trace of the positive
    semidefinite normalized pure state, and the partial trace preserves the trace, which
    the normalization fixes to one.
\end{proof}

\begin{proposition}[Pure-state area-law upper bound]\label{prop:pure_area_upper}
    \lean{MPSTensor.pureBlockEntropy_le}
    \leanok
    \uses{lem:operator_schmidt_rank, thm:entropy_le_log_rank,
        lem:reduced_pure_block_state_density}
    Let $A$ be an MPS tensor of bond dimension $D\ge 1$ with $V^{(N)}(A)\neq 0$.
    For every block length $L\le N$,
    \[
        S_L^{(N)}(A) \le 2\log D,
    \]
    uniformly in $L$ and $N$. This is the area law of
    \cite[Section~3, line~599]{Cirac2016MPDO_arXiv}: the block entropy is bounded
    by a constant independent of the block size.
\end{proposition}
\begin{proof}\leanok
    \uses{lem:reduced_pure_block_state_density}
    The reduced state $\rho_L^{(N)}(A)$ is a density matrix
    (Lemma~\ref{lem:reduced_pure_block_state_density}) of rank at most $D^2$ by
    Lemma~\ref{lem:operator_schmidt_rank}. By the rank bound on the entropy
    (Theorem~\ref{thm:entropy_le_log_rank}),
    $S_L^{(N)}(A) \le \log(D^2) = 2\log D$.
\end{proof}

\begin{lemma}[Nonnegativity of the pure-state block entropy]
    \label{lem:pure_block_entropy_nonneg}
    \lean{MPSTensor.pureBlockEntropy_nonneg}
    \leanok
    \uses{def:pure_sal, thm:entropy_nonneg, lem:reduced_pure_block_state_density}
    For an MPS tensor $A$ with $V^{(N)}(A)\neq 0$ and every block length $L\le N$,
    the pure-state block entropy is nonnegative: $0 \le S_L^{(N)}(A)$.
\end{lemma}
\begin{proof}\leanok
    \uses{lem:reduced_pure_block_state_density}
    The reduced state $\rho_L^{(N)}(A)$ is a density matrix
    (Lemma~\ref{lem:reduced_pure_block_state_density}), and the von Neumann
    entropy of a density matrix is nonnegative.
\end{proof}

\begin{lemma}[Full-block pure-state entropy vanishes]
    \label{lem:pure_entropy_full}
    \lean{pureBlockEntropy_full}
    \leanok
    \uses{def:pure_sal}
    For an MPS tensor $A$ with $V^{(N)}(A)\neq 0$, the entropy of the whole chain
    vanishes:
    \[
        S_N^{(N)}(A) = 0 .
    \]
\end{lemma}
\begin{proof}\leanok
    The normalized operator is the rank-one pure state
    \[
        \frac{\ket{V^{(N)}(A)}\bra{V^{(N)}(A)}}
        {\langle V^{(N)}(A), V^{(N)}(A)\rangle} ,
    \]
    after reindexing the full block. A pure state has von Neumann entropy zero.
\end{proof}

\begin{lemma}[Empty-block pure-state entropy vanishes]
    \label{lem:pure_block_entropy_empty}
    \lean{MPSTensor.pureBlockEntropy_empty}
    \leanok
    \uses{def:pure_sal, lem:pure_entropy_full, lem:schmidt_symmetry}
    For an MPS tensor $A$ with $V^{(N)}(A)\neq 0$, the pure-state block entropy of the
    empty block vanishes: $S_0^{(N)}(A) = 0$.
\end{lemma}
\begin{proof}\leanok
    \uses{lem:pure_entropy_full, lem:schmidt_symmetry}
    By the Schmidt symmetry $S_0^{(N)}(A) = S_{N}^{(N)}(A)$, and the entropy of the full
    system vanishes for a pure state (Lemma~\ref{lem:pure_entropy_full}).
\end{proof}

\begin{lemma}[Saturation telescopes to a constant block entropy]
    \label{lem:pure_sal_const}
    \lean{MPSTensor.pureBlockEntropy_eq_of_isSAL}
    \leanok
    \uses{def:pure_sal}
    If an MPS tensor $A$ saturates the area law, then all block entropies in the
    range $1\le L\le\lfloor N/2\rfloor$ coincide:
    \[
        S_L^{(N)}(A) = S_{L'}^{(N)}(A)
        \qquad
        \text{for all } 1\le L, L'\le \lfloor N/2\rfloor .
    \]
\end{lemma}
\begin{proof}\leanok
    \uses{def:pure_sal}
    The defining condition supplies the single-step equalities
    $S_m^{(N)}(A) = S_{m+1}^{(N)}(A)$ for $1\le m<\lfloor N/2\rfloor$. Climbing one
    step at a time from the smaller of $L$ and $L'$ up to the larger composes these
    into the stated equality.
\end{proof}

\begin{lemma}[Pure-state mutual information is twice the block entropy]
    \label{lem:mutual_info_doubled}
    \lean{mutualInfoChain_doubledTensor}
    \leanok
    \uses{lem:pure_entropy_full, lem:schmidt_symmetry, lem:blockEntropy_doubledTensor}
    For an MPS tensor $A$ with $V^{(N)}(A)\neq 0$, the mutual information of the
    purification MPDO equals twice the pure-state block entropy:
    \[
        I_L = 2\,S_L^{(N)}(A).
    \]
\end{lemma}
\begin{proof}\leanok
    By definition $I_L = S_L + S_{N-L} - S_N$ for the generated state. The global
    entropy of the pure state vanishes (Lemma~\ref{lem:pure_entropy_full}), and
    the Schmidt symmetry (Lemma~\ref{lem:schmidt_symmetry}) gives
    $S_{N-L} = S_L$, so $I_L = 2\,S_L$.
\end{proof}

\begin{proposition}[Pure-state mutual-information bound]\label{prop:pure_mutual_info_bound}
    \lean{mutualInfoChain_doubledTensor_le}
    \leanok
    \uses{prop:pure_area_upper, lem:mutual_info_doubled}
    Let $A$ be an MPS tensor of bond dimension $D\ge 1$ with $V^{(N)}(A)\neq 0$.
    For the purification MPDO of $A$, whose generated operator is the pure state
    $|V^{(N)}(A)\rangle\langle V^{(N)}(A)|$, the mutual information between a block
    of $L$ spins and the rest satisfies
    \[
        I_L \le 4\log D,
    \]
    uniformly in $L$ and $N$.
\end{proposition}
\begin{proof}\leanok
    For a pure state $I_L = 2\,S_L^{(N)}(A)$, since the global entropy $S_N$
    vanishes and a block and its complement share the same Schmidt spectrum. The
    block-entropy bound $S_L^{(N)}(A) \le 2\log D$
    (Proposition~\ref{prop:pure_area_upper}) then gives $I_L \le 4\log D$.
\end{proof}
